
égalités 2.4.2 sont des polynômes en \(X,Y\), à coefficients rationnels, et il suffit donc, pour montrer leur égalité, de montrer qu'ils prennent même valeur pour toutes valeurs entières de \(X,Y\). Mais les relations obtenues par substitution de valeurs entières à \(X,Y\) sont celles qui expriment que le pré- -anneau \(Z\) (2.5.a)) est un -anneau; elles sont donc satisfaites.

\textbf{Définition 2.7.} \uline{On appelle} anneau binômial \uline{un anneau vérifiant les conditions} 2.6.i) \uline{et} ii), \uline{muni de la structure de} \(\lambda\)-\uline{anneau définie par} 2.6.1.

\textbf{Exemple 2.8.} Soient \(X\) un site, \(K\) un anneau binômial, \(\underline{K}\) le faisceau associé au préfaisceau constant de valeur \(K\) sur \(X\) ; alors \(H^0(X,\underline{K})\) est un anneau binômial.

En effet, les conditions 2.6 se conservent par passage à la limite inductive.

En particulier, \(H^0(X,\underline{\mathbb Z})\) est un anneau binômial.

2.9. Soient \(K\) un anneau binômial, \(A\) une \(K\)-algèbre. Pour toute série \(\varphi\in\Ghat(A)\), et tout \(a\in K\), on peut définir une série \(\varphi^a\) de la façon usuelle : on écrit \(\varphi=1+\psi\), et on pose :
\begin{equation}\tag{2.9.1}
  \varphi^a=(1+\psi)^a
  =
  \sum_{n=0}^{\infty}\binom{a}{n}\psi^n
  \left(=\exp(a\Log\varphi)\right),
\end{equation}
la dernière égalité ne pouvant être écrite que modulo torsion. Ces formules sont obtenues par substitution à partir des formules analogues pour \(\varphi=1+t\) ; les relations suivantes, propriétés classiques de l'exponentielle, restent donc encore vraies pour \(a,b\in K\), \(\varphi,\psi\in\Ghat(A)\) :
\begin{equation}\tag{2.9.2}
  \varphi^a\varphi^b=\varphi^{a+b};\qquad
  \varphi^{ab}=(\varphi^a)^b;\qquad
  \varphi^a\psi^a=(\varphi\psi)^a .
\end{equation}



% Exposé V continuation macros
\providecommand{\Ghat}{\widehat G}
\providecommand{\Gg}{\widehat G_g}
\providecommand{\Gcirc}{\widehat G_{\circ}}
\providecommand{\Gnabla}{\widehat G_{\nabla}}
\providecommand{\Grhat}{\widehat{\operatorname{Gr}}}
\providecommand{\Log}{\operatorname{Log}}
\providecommand{\Fil}{\operatorname{Fil}}
\providecommand{\eps}{\varepsilon}
\renewcommand{\id}{\operatorname{id}}
\providecommand{\Sym}{\operatorname{Sym}}
\providecommand{\Mon}{\operatorname{Mon}}


